\section{Discussion}
\label{sec:discussion}

\para{What worked?} The strongest positive result is not accuracy, but controllability. A simple JSON speaker-ID prompt produced valid outputs in all 60 LLM calls and preserved transcript text by construction. This matters because prior LLM post-processing can degrade transcripts when the model is allowed to rewrite words~\cite{wang2024diarizationlm}. In this pilot, the LLM correction layer is safe in the narrow sense that it does not alter ASR text.

\para{What failed?} The LLM rarely changed labels in four-speaker \amisdm windows. Many local speaker IDs are anonymized, and short meeting snippets do not provide stable lexical identities. Without acoustic confidence, embeddings, names, roles, or longer history, a text-only LLM has little evidence for deciding that speaker \textsc{S1} should become speaker \textsc{S3}. This explains why both LLM conditions exactly match \noisylabels on \amisdm. The model's conservatism is preferable to hallucinated relabeling, but it limits correction value.

\para{Why did profiles not help?} We expected profile-conditioned relabeling to outperform no-profile relabeling because profile turns could anchor each speaker. The result does not support that hypothesis. The profiles were short, text-only, and unscored; they did not provide robust identity evidence in meeting snippets. In \simsamu, the profile condition improves over \noisylabels but underperforms \llmnoprofile. This suggests that weak profiles can add prompt complexity without adding enough speaker evidence.

\para{Why did continuity smoothing hurt?} \smooth assumes that isolated speaker changes are likely errors. That assumption fails in dialogue where true turns alternate quickly. The method increases mean WDER from 0.390 to 0.442, and it is especially poor in \simsamu low-noise cases. This failure is useful: a transcript correction layer needs discourse-aware evidence, not only local temporal smoothness.

\para{Limitations.} This is a pilot, not a benchmark claim. The experiment has five windows and 30 noisy cases. The \amisdm shard lacks manual word transcripts, so ASR text is a pseudo-reference for words and the evaluation isolates speaker-label correction rather than ASR word accuracy. Noise injection is controlled and seeded, but it simulates diarization errors rather than using a full WhisperX plus pyannote cascade. The speaker profiles are short and text-only. Finally, our cpWER-style score is a local exact-permutation implementation for small speaker counts, not an official MeetEval run.

\para{Broader implication.} The result points to a practical design rule: use prompted LLMs as constrained, transcript-preserving relabelers only when they receive stronger speaker evidence. Acoustic confidence scores, speaker embeddings, enrollment profiles, longer conversational history, or fine-tuned DiarizationLM-style models are likely necessary for robust many-speaker meeting identification.
